%-----------------------------------------------------------------------------
%
%               Template for sigplanconf LaTeX Class
%
% Name:         sigplanconf-template.tex
%
% Purpose:      A template for sigplanconf.cls, which is a LaTeX 2e class
%               file for SIGPLAN conference proceedings.
%
% Guide:        Refer to "Author's Guide to the ACM SIGPLAN Class,"
%               sigplanconf-guide.pdf
%
% Author:       Paul C. Anagnostopoulos
%               Windfall Software
%               978 371-2316
%               paul@windfall.com
%
% Created:      15 February 2005 
%
%-----------------------------------------------------------------------------


\documentclass[10pt,draft]{sigplanconf}

% The following \documentclass options may be useful:

% preprint      Remove this option only once the paper is in final form.
% 10pt          To set in 10-point type instead of 9-point.
% 11pt          To set in 11-point type instead of 9-point.
% authoryear    To obtain author/year citation style instead of numeric.

\usepackage{amsmath}
\usepackage[obeyDraft]{todonotes}

\begin{document}
\presetkeys{todonotes}{inline}{}

\special{papersize=8.5in,11in}
\setlength{\pdfpageheight}{\paperheight}
\setlength{\pdfpagewidth}{\paperwidth}

\conferenceinfo{DLS '15}{October 25--30, 2015, Pittsburgh, PA, USA} 
\copyrightyear{2025} 
\copyrightdata{978-1-nnnn-nnnn-n/yy/mm} 
\doi{nnnnnnn.nnnnnnn}

% Uncomment one of the following two, if you are not going for the 
% traditional copyright transfer agreement.

\exclusivelicense                % ACM gets exclusive license to publish, 
                                  % you retain copyright

%\permissiontopublish             % ACM gets nonexclusive license to publish
                                  % (paid open-access papers, 
                                  % short abstracts)

\titlebanner{banner above paper title}        % These are ignored unless
\preprintfooter{short description of paper}   % 'preprint' option specified.

\title{DLS Paper about RStrategies}
\subtitle{Subtitle Text, if any}

\authorinfo{Anton Gulenko\and Tobias Pape\and Tim~Felgentreff\and Robert Hirschfeld}
           {Hasso Plattner Institute, University of Potsdam}
           {\{firstname.lastname\}@hpi.uni-potsdam.de}
\authorinfo{Carl Friedrich Bolz}
           {King's College, London}
           {cfbolz@gmx.de}

\maketitle

\begin{abstract}
This is the text of the abstract.
\end{abstract}

\category{CR-number}{subcategory}{third-level}

% general terms are not compulsory anymore, 
% you may leave them out
\terms
term1, term2

\keywords
keyword1, keyword2

\section{Introduction}

\todo{Contributions ... need to think about those, below just suggestions}
\begin{itemize}
\item An application of strategies to object storage in RSqueak
\item A generic framework to add strategies to other PyPy-based languages
\item An evaluation of the strategies approach for a wider range of languages
\end{itemize}

\section{Background: Storage Strategies in PyPy}

\section{Language Independent Storage Strategies}

\subsection{The rstrategies Framework}
\todo{Framework description}

\subsection{Storage Strategies for Objects}
\todo{RSqueak Implementation}
\todo{which strategies are most important? why?}

\subsubsection{Single Value Strategies}
\todo{Int, Float, Nil}

\subsubsection{Mixed Value Strategies}
\todo{FloatAndInt, NilAndInt, NilAndFloat?}

\subsection{Run-Length Strategies}
\todo{Ranges, length-of-nil, ...?  Do we have this? Should we have this?
Otherwise, we can just move it to future work...?}

\subsection{Supporting Exotic Objects}
\todo{ContextParts, Classes, MethodDicts}

\section{Evaluation}

\subsection{Performance}
\todo{Does the performance benefit reported in the PyPy paper carry over?}
\begin{table}
\label{fig:performance}
\caption{Table shows...}
\missingfigure{\Huge Need the figure here!}
\end{table}

\subsection{Memory}
\todo{Does the memory benefit carry over?}
\begin{figure}
\label{fig:memory}
\caption{Memory overhead...}
\missingfigure{\Huge Need the figure here!}
\end{figure}

\subsection{Generality}
\todo{Can the framework be applied to other languages? Topaz, Pycket}

\subsection{Transitions}
\todo{Are the strategies we identified as most important really that?}
\begin{figure}
\label{fig:transitions}
\caption{Transitions...}
\missingfigure{\Huge Need the figure here!}
\end{figure}

\section{Related Work}

\todo{Strategies with Tracing JIT, PyPy}
\todo{Strategies with Method JIT (Truffle, V8)}
\todo{Pointer tagging}

\section{Discussion and Conclusion}
\todo{which further strategies do we need?}
\todo{future work}

\bibliography{dls}
\bibliographystyle{abbrvnat}


\end{document}

%                       Revision History
%                       -------- -------
%  Date         Person  Ver.    Change
%  ----         ------  ----    ------

%  2013.06.29   TU      0.1--4  comments on permission/copyright notices